\documentclass[10pt,a4paper,twocolumn]{article}
\usepackage{subfiles}

% Shared preamble (packages, colors, listings, theorem environments)
% ==============================================================================
% banq-preamble.tex
% Shared LaTeX preamble for all XPower Banq papers and the unified bundle.
% Loaded via % ==============================================================================
% banq-preamble.tex
% Shared LaTeX preamble for all XPower Banq papers and the unified bundle.
% Loaded via % ==============================================================================
% banq-preamble.tex
% Shared LaTeX preamble for all XPower Banq papers and the unified bundle.
% Loaded via \input{../shared/banq-preamble} from each child paper and from
% the master bundle root. Identical to the original per-paper preambles.
% ==============================================================================

\usepackage[utf8]{inputenc}
\usepackage[T1]{fontenc}
\usepackage{newunicodechar}
\usepackage{libertinus}
\usepackage[margin=2cm]{geometry}
\usepackage{amsmath,amssymb,amsthm}
\usepackage{graphicx}
\usepackage{xcolor}
\usepackage[hyphens]{url}
\usepackage{hyperref}
\usepackage{booktabs}
\usepackage{tabularx}
\usepackage{multirow}
\usepackage{enumitem}
\usepackage{tikz}
\usepackage{pgfplots}
\usepackage{listings}
\usepackage{caption}
\usepackage{subcaption}
\usepackage{fancyhdr}
\usepackage{float}
\usepackage{placeins}  % For \FloatBarrier
\usepackage{tocloft}   % Adjustable TOC number-box widths
\usepackage{multicol}  % Multi-column layout (used for the bundle's master TOC)
\usepackage{twemojis}  % Twemoji glyphs (pdflatex-compatible color emoji)

% Float parameters to prevent overlapping
\setcounter{topnumber}{2}
\setcounter{bottomnumber}{2}
\setcounter{totalnumber}{4}
\renewcommand{\topfraction}{0.85}
\renewcommand{\bottomfraction}{0.85}
\renewcommand{\textfraction}{0.15}
\renewcommand{\floatpagefraction}{0.7}

% Prevent line breaks inside inline math (fixes broken superscripts/sqrt)
\relpenalty=10000
\binoppenalty=10000

\pgfplotsset{compat=1.18}
\usetikzlibrary{shapes,arrows,positioning,calc,patterns}

% ==============================================================================
% Colors and Styling
% ==============================================================================
\definecolor{codeblue}{RGB}{0,102,204}
\definecolor{codepurple}{RGB}{128,0,128}
\definecolor{codegreen}{RGB}{0,128,0}
\definecolor{codegray}{RGB}{128,128,128}
\definecolor{primaryblue}{RGB}{41,98,255}
\definecolor{warningred}{RGB}{220,53,69}
\definecolor{successgreen}{RGB}{40,167,69}

\hypersetup{
    colorlinks=true,
    linkcolor=primaryblue,
    citecolor=primaryblue,
    urlcolor=primaryblue
}

\lstset{
    basicstyle=\ttfamily\scriptsize,
    keywordstyle=\color{codeblue}\bfseries,
    commentstyle=\color{codegreen},
    stringstyle=\color{codepurple},
    numbers=left,
    numberstyle=\tiny\color{codegray},
    numbersep=4pt,
    xleftmargin=2em,
    framexleftmargin=2em,
    breaklines=true,
    frame=single,
    captionpos=b
}
\lstdefinelanguage{Solidity}{
    keywords={function, returns, return, if, else, for, while, do, break, continue, throw, require, revert, assert, public, private, internal, external, view, pure, payable, memory, storage, calldata, constant, immutable, virtual, override, modifier, event, emit, struct, enum, mapping, contract, interface, library, abstract, is, new, delete, true, false, uint256, uint128, uint64, uint32, uint16, uint8, uint240, uint224, int256, int128, int64, int32, int8, address, bool, bytes, bytes32, string, unchecked},
    sensitive=true,
    morecomment=[l]{//},
    morecomment=[s]{/*}{*/},
    morestring=[b]",
    morestring=[b]'
}

% ==============================================================================
% Theorem Environments
% ==============================================================================
\theoremstyle{definition}
\newtheorem{definition}{Definition}[section]
\newtheorem{property}{Property}[section]
\theoremstyle{plain}
\newtheorem{theorem}{Theorem}[section]
\newtheorem{lemma}[theorem]{Lemma}
\newtheorem{corollary}[theorem]{Corollary}
\theoremstyle{remark}
\newtheorem{remark}{Remark}[section]
 from each child paper and from
% the master bundle root. Identical to the original per-paper preambles.
% ==============================================================================

\usepackage[utf8]{inputenc}
\usepackage[T1]{fontenc}
\usepackage{newunicodechar}
\usepackage{libertinus}
\usepackage[margin=2cm]{geometry}
\usepackage{amsmath,amssymb,amsthm}
\usepackage{graphicx}
\usepackage{xcolor}
\usepackage[hyphens]{url}
\usepackage{hyperref}
\usepackage{booktabs}
\usepackage{tabularx}
\usepackage{multirow}
\usepackage{enumitem}
\usepackage{tikz}
\usepackage{pgfplots}
\usepackage{listings}
\usepackage{caption}
\usepackage{subcaption}
\usepackage{fancyhdr}
\usepackage{float}
\usepackage{placeins}  % For \FloatBarrier
\usepackage{tocloft}   % Adjustable TOC number-box widths
\usepackage{multicol}  % Multi-column layout (used for the bundle's master TOC)
\usepackage{twemojis}  % Twemoji glyphs (pdflatex-compatible color emoji)

% Float parameters to prevent overlapping
\setcounter{topnumber}{2}
\setcounter{bottomnumber}{2}
\setcounter{totalnumber}{4}
\renewcommand{\topfraction}{0.85}
\renewcommand{\bottomfraction}{0.85}
\renewcommand{\textfraction}{0.15}
\renewcommand{\floatpagefraction}{0.7}

% Prevent line breaks inside inline math (fixes broken superscripts/sqrt)
\relpenalty=10000
\binoppenalty=10000

\pgfplotsset{compat=1.18}
\usetikzlibrary{shapes,arrows,positioning,calc,patterns}

% ==============================================================================
% Colors and Styling
% ==============================================================================
\definecolor{codeblue}{RGB}{0,102,204}
\definecolor{codepurple}{RGB}{128,0,128}
\definecolor{codegreen}{RGB}{0,128,0}
\definecolor{codegray}{RGB}{128,128,128}
\definecolor{primaryblue}{RGB}{41,98,255}
\definecolor{warningred}{RGB}{220,53,69}
\definecolor{successgreen}{RGB}{40,167,69}

\hypersetup{
    colorlinks=true,
    linkcolor=primaryblue,
    citecolor=primaryblue,
    urlcolor=primaryblue
}

\lstset{
    basicstyle=\ttfamily\scriptsize,
    keywordstyle=\color{codeblue}\bfseries,
    commentstyle=\color{codegreen},
    stringstyle=\color{codepurple},
    numbers=left,
    numberstyle=\tiny\color{codegray},
    numbersep=4pt,
    xleftmargin=2em,
    framexleftmargin=2em,
    breaklines=true,
    frame=single,
    captionpos=b
}
\lstdefinelanguage{Solidity}{
    keywords={function, returns, return, if, else, for, while, do, break, continue, throw, require, revert, assert, public, private, internal, external, view, pure, payable, memory, storage, calldata, constant, immutable, virtual, override, modifier, event, emit, struct, enum, mapping, contract, interface, library, abstract, is, new, delete, true, false, uint256, uint128, uint64, uint32, uint16, uint8, uint240, uint224, int256, int128, int64, int32, int8, address, bool, bytes, bytes32, string, unchecked},
    sensitive=true,
    morecomment=[l]{//},
    morecomment=[s]{/*}{*/},
    morestring=[b]",
    morestring=[b]'
}

% ==============================================================================
% Theorem Environments
% ==============================================================================
\theoremstyle{definition}
\newtheorem{definition}{Definition}[section]
\newtheorem{property}{Property}[section]
\theoremstyle{plain}
\newtheorem{theorem}{Theorem}[section]
\newtheorem{lemma}[theorem]{Lemma}
\newtheorem{corollary}[theorem]{Corollary}
\theoremstyle{remark}
\newtheorem{remark}{Remark}[section]
 from each child paper and from
% the master bundle root. Identical to the original per-paper preambles.
% ==============================================================================

\usepackage[utf8]{inputenc}
\usepackage[T1]{fontenc}
\usepackage{newunicodechar}
\usepackage{libertinus}
\usepackage[margin=2cm]{geometry}
\usepackage{amsmath,amssymb,amsthm}
\usepackage{graphicx}
\usepackage{xcolor}
\usepackage[hyphens]{url}
\usepackage{hyperref}
\usepackage{booktabs}
\usepackage{tabularx}
\usepackage{multirow}
\usepackage{enumitem}
\usepackage{tikz}
\usepackage{pgfplots}
\usepackage{listings}
\usepackage{caption}
\usepackage{subcaption}
\usepackage{fancyhdr}
\usepackage{float}
\usepackage{placeins}  % For \FloatBarrier
\usepackage{tocloft}   % Adjustable TOC number-box widths
\usepackage{multicol}  % Multi-column layout (used for the bundle's master TOC)
\usepackage{twemojis}  % Twemoji glyphs (pdflatex-compatible color emoji)

% Float parameters to prevent overlapping
\setcounter{topnumber}{2}
\setcounter{bottomnumber}{2}
\setcounter{totalnumber}{4}
\renewcommand{\topfraction}{0.85}
\renewcommand{\bottomfraction}{0.85}
\renewcommand{\textfraction}{0.15}
\renewcommand{\floatpagefraction}{0.7}

% Prevent line breaks inside inline math (fixes broken superscripts/sqrt)
\relpenalty=10000
\binoppenalty=10000

\pgfplotsset{compat=1.18}
\usetikzlibrary{shapes,arrows,positioning,calc,patterns}

% ==============================================================================
% Colors and Styling
% ==============================================================================
\definecolor{codeblue}{RGB}{0,102,204}
\definecolor{codepurple}{RGB}{128,0,128}
\definecolor{codegreen}{RGB}{0,128,0}
\definecolor{codegray}{RGB}{128,128,128}
\definecolor{primaryblue}{RGB}{41,98,255}
\definecolor{warningred}{RGB}{220,53,69}
\definecolor{successgreen}{RGB}{40,167,69}

\hypersetup{
    colorlinks=true,
    linkcolor=primaryblue,
    citecolor=primaryblue,
    urlcolor=primaryblue
}

\lstset{
    basicstyle=\ttfamily\scriptsize,
    keywordstyle=\color{codeblue}\bfseries,
    commentstyle=\color{codegreen},
    stringstyle=\color{codepurple},
    numbers=left,
    numberstyle=\tiny\color{codegray},
    numbersep=4pt,
    xleftmargin=2em,
    framexleftmargin=2em,
    breaklines=true,
    frame=single,
    captionpos=b
}
\lstdefinelanguage{Solidity}{
    keywords={function, returns, return, if, else, for, while, do, break, continue, throw, require, revert, assert, public, private, internal, external, view, pure, payable, memory, storage, calldata, constant, immutable, virtual, override, modifier, event, emit, struct, enum, mapping, contract, interface, library, abstract, is, new, delete, true, false, uint256, uint128, uint64, uint32, uint16, uint8, uint240, uint224, int256, int128, int64, int32, int8, address, bool, bytes, bytes32, string, unchecked},
    sensitive=true,
    morecomment=[l]{//},
    morecomment=[s]{/*}{*/},
    morestring=[b]",
    morestring=[b]'
}

% ==============================================================================
% Theorem Environments
% ==============================================================================
\theoremstyle{definition}
\newtheorem{definition}{Definition}[section]
\newtheorem{property}{Property}[section]
\theoremstyle{plain}
\newtheorem{theorem}{Theorem}[section]
\newtheorem{lemma}[theorem]{Lemma}
\newtheorem{corollary}[theorem]{Corollary}
\theoremstyle{remark}
\newtheorem{remark}{Remark}[section]

% ==============================================================================
% banq-meta.tex
% Shared author/version metadata for all XPower Banq papers.
% ==============================================================================

\newcommand{\banqAuthor}{\textsc{Karun The Ritch}}
\newcommand{\banqAuthorEmail}{ktr@xpowerbanq.com}
\newcommand{\banqAuthorURL}{https://www.xpowerbanq.com}
\newcommand{\banqAuthorContact}{%
    \href{mailto:\banqAuthorEmail}{\texttt{\banqAuthorEmail}}\quad\url{\banqAuthorURL}%
}
\newcommand{\banqArxivId}{preprint:2605~~[q-fin.CP]}
\newcommand{\banqArxivDate}{1 May 2026}

% ==============================================================================
% banq-bundle.tex
% Defines the \ifbundle conditional that distinguishes a standalone child build
% from the unified bundle build, and the \xref macro that produces a real
% internal hyperlink in bundle mode and a citation-with-text in standalone mode.
%
% Standalone child papers leave \bundlefalse (the default).
% The master bundle root sets \bundletrue before \begin{document}.
% ==============================================================================

\newif\ifbundle
\bundlefalse

% \xref{<bibkey>}{<label>}{<text>}
% - In bundle mode: produces a clickable in-document hyperlink to the label.
% - In standalone mode: prints the text followed by a citation to the bibkey.
% Note: defined after \bundletrue/\bundlefalse so it picks up the right branch.
% Children must call this from inside \AtBeginDocument to ensure \ifbundle is
% finalized; alternatively the macro can test \ifbundle at call site.
\newcommand{\xref}[3]{%
  \ifbundle
    \hyperref[#2]{#3}%
  \else
    #3~\cite{#1}%
  \fi
}

% \xrefshort{<bibkey>}{<label>}{<text>}
% Same as \xref but suppresses the citation entirely in standalone mode (used
% when the surrounding prose already cites the source).
\newcommand{\xrefshort}[3]{%
  \ifbundle
    \hyperref[#2]{#3}%
  \else
    #3%
  \fi
}

% \banqEnableBundleNumbering
% Resets the section counter at every \part and prefixes \thesection with the
% part's Roman numeral, so headings read "I.1 Introduction", "II.1 Introduction",
% etc. Called explicitly from the master root after \bundletrue (see
% 000.banq-all/banq.tex). Not called by standalone children.
%
% Parts that contain MULTIPLE sub-papers must additionally call \banqSubpart
% (defined below) before each child \subfile, which appends a per-paper letter
% suffix (a, b, ...) to the part numeral and resets the section counter so each
% sub-paper restarts from 1. Without this, sub-papers within the same part
% share a single section sequence, which causes the TOC to conflate identically
% titled sections (e.g. two "II.1 Introduction" entries for banq-lck/banq-log).
%
% The outer \makeatletter/\makeatother brackets here are critical: the macro
% body contains \@addtoreset (an @ control sequence) which must be parseable
% as a single token at \newcommand-definition time, not at call time.
\newcommand{\banqSubpartLetter}{}
% Boolean tracking whether we are currently inside a sub-paper of a shared part
% (i.e. \banqSubpart has been called since the last \part). Used by \banqPaper
% to decide whether to demote section bookmark levels. \ifx-comparing
% \banqSubpartLetter against \@empty doesn't work because \newcommand produces
% \long macros while \@empty is short — \ifx always returns false.
\newif\ifbanqSharedPart
\makeatletter
\newcommand{\banqEnableBundleNumbering}{%
  \@addtoreset{section}{part}%
  \renewcommand{\thesection}{\Roman{part}\banqSubpartLetter.\,\arabic{section}}%
  % Make hyperref anchor names unique across sub-papers that share a \part.
  % Without this, banq-lck and banq-log both emit section.2.1, section.2.2, …
  % which corrupts the PDF outline tree (external viewers truncate at Part II).
  \renewcommand{\theHsection}{\Roman{part}\banqSubpartLetter.\arabic{section}}%
  \renewcommand{\theHsubsection}{\theHsection.\arabic{subsection}}%
  \renewcommand{\theHsubsubsection}{\theHsubsection.\arabic{subsubsection}}%
  \renewcommand{\theHparagraph}{\theHsubsubsection.\arabic{paragraph}}%
  % Hook \part so the subpart-letter and shared-part flag both reset when a new
  % part begins; otherwise state set in part N would leak into part N+1.
  \let\banqOriginalPart\part
  \renewcommand{\part}{%
    \renewcommand{\banqSubpartLetter}{}%
    \banqSharedPartfalse
    \banqOriginalPart
  }%
  % The "I.\,10" prefix is wider than the default \numberline box. Widen the
  % TOC number-box widths so two-digit section numbers don't overlap titles.
  \setlength{\cftsecnumwidth}{3.5em}%
  \setlength{\cftsubsecnumwidth}{4.2em}%
  \setlength{\cftsubsubsecnumwidth}{5.2em}%
}
\makeatother

% \banqSubpart{<letter>}
% Marks the start of a sub-paper inside a part that contains multiple papers.
% Sets the per-paper letter suffix appended to the part numeral in \thesection
% (so headings read "IIa.1 Introduction", "IIb.1 Introduction", etc.) and
% resets the section counter so the sub-paper's own numbering restarts at 1.
% Pass an empty argument (\banqSubpart{}) to clear the suffix mid-part if
% needed; the suffix also auto-clears at every \part boundary (see hook above).
\newcommand{\banqSubpart}[1]{%
  \renewcommand{\banqSubpartLetter}{#1}%
  \setcounter{section}{0}%
  \banqSharedParttrue
}

% \banq@RestoreLevels / \banq@DemoteLevels
% Adjust hyperref's per-sectioning bookmark levels. When two sub-papers share a
% \part (banq-lck "IIa" and banq-log "IIb" inside Part~II), each paper's
% sections should nest under its own divider bookmark in the PDF outline,
% rather than appearing as siblings of it under the part. Demoting
% \toclevel@section from 1 to 2 (and the deeper levels accordingly) achieves
% that nesting; \banqPaper calls Restore before emitting its own divider so the
% divider always lands at level 1, then calls Demote afterwards (only when in a
% shared part) so subsequent sections sit at level 2 = child of the divider.
% Affects only the PDF outline depths; the printed TOC layout is unchanged.
\makeatletter
\newcommand{\banq@RestoreLevels}{%
  \renewcommand{\toclevel@section}{1}%
  \renewcommand{\toclevel@subsection}{2}%
  \renewcommand{\toclevel@subsubsection}{3}%
  \renewcommand{\toclevel@paragraph}{4}%
}
\newcommand{\banq@DemoteLevels}{%
  \renewcommand{\toclevel@section}{2}%
  \renewcommand{\toclevel@subsection}{3}%
  \renewcommand{\toclevel@subsubsection}{4}%
  \renewcommand{\toclevel@paragraph}{5}%
}
\makeatother

% Storage for the current part's twemoji name. \banqPartWithEmoji sets it
% (overwritten at each part boundary). \banqPaper reads it without clearing,
% so multiple sub-papers within the same part (e.g. banq-lck and banq-log
% in Part~II) share the same emoji on their respective cover pages.
\makeatletter
\def\banq@partemoji{}
\newif\ifbanq@hasemoji
\banq@hasemojifalse
% \banqSetSubpartEmoji{<twemoji-name>}
% Override the part's default emoji for the next sub-paper cover. Used inside
% parts that contain multiple papers when each sub-paper deserves its own
% cover icon (e.g. banq-lck takes "lock" while banq-log keeps the part-level
% "gear"). Encapsulates the @-internal \gdef so the master document body need
% not enter \makeatletter scope.
\newcommand{\banqSetSubpartEmoji}[1]{\gdef\banq@partemoji{#1}}
\makeatother

% \banqForceRecto
% Page-break helper that flushes outstanding output and ensures the next
% content-bearing page lands on an odd (recto) page in duplex bindings.
% The article class is loaded without `twoside` so \cleardoublepage degenerates
% to \clearpage and will not pad; we check the page counter explicitly and
% insert a blank filler when the upcoming page is even. \value{page} reflects
% the next page number after \clearpage, before any content has been emitted
% on it. Call sites: front-matter (TOC start, body start) and \banqPaper.
\newcommand{\banqForceRecto}{%
  \clearpage
  \ifodd\value{page}\else\hbox{}\thispagestyle{empty}\newpage\fi
}

% \banqEnsureEvenPageCount
% Final-pass parity helper: guarantees the document ends with an even total
% page count, so duplex-bound copies have a back cover on the verso side and
% any subsequent insert (e.g. a colophon, a bound-in errata sheet) starts on
% a recto. After \clearpage, \value{page} holds the *next* page number, so
% the current page count equals \value{page}-1; the count is even iff
% \value{page} is odd, which is exactly the invariant \banqForceRecto already
% enforces. The two operations are semantically distinct (recto-start vs.
% even-total) but compute identical pad logic, so we alias them. Call once,
% immediately before \end{document}, in the bundle root.
\let\banqEnsureEvenPageCount\banqForceRecto

% \banqPartWithEmoji{<title>}{<twemoji-name>}
% Bumps the part counter, stashes the title for the part-level TOC entry,
% and stashes the emoji for the cover page(s) that follow (see \banqPaper).
% No physical divider page is emitted — the cover that immediately follows
% is the visible boundary between parts. The bundle's section-numbering
% scheme (I.1, IIa.1, IIb.1, III.1, ...) continues to work because
% \refstepcounter drives the same counter \part would have, and
% \@addtoreset (set up by \banqEnableBundleNumbering) still resets section
% at the part boundary.
%
% The TOC entry itself is written by \banqPaper at the start of the cover
% page rather than here, so the deferred page-number capture lands on the
% cover (the visible part boundary) rather than on whatever page the build
% happens to be on between subfiles.
\makeatletter
\def\banq@parttitle{}
\newif\ifbanq@parttocpending
\banq@parttocpendingfalse
\newcommand{\banqPartWithEmoji}[2]{%
  \refstepcounter{part}%
  \renewcommand{\banqSubpartLetter}{}%
  \banqSharedPartfalse
  \gdef\banq@partemoji{#2}%
  \global\banq@hasemojitrue
  \gdef\banq@parttitle{#1}%
  \global\banq@parttocpendingtrue
}
\makeatother

% \banqAbstract{<abstract text>}{<keywords>}
% Single source of truth for each child's abstract + keywords list, with
% mode-aware rendering:
%   - Standalone: emits a normal \begin{abstract}...\end{abstract} block
%     followed by a "Keywords:" line (the previous \ifbundle\else ... \fi
%     boilerplate each child used to inline).
%   - Bundle: stashes the content into module-level token lists and sets
%     \ifbanq@hasabstract; \banqPaper then renders it on the verso page that
%     follows the cover (the page that was previously left blank).
% Defined as \long so abstracts containing paragraph breaks are accepted,
% though in practice all six abstracts are single-paragraph.
\makeatletter
\def\banq@abstracttext{}
\def\banq@keywordstext{}
\newif\ifbanq@hasabstract
\banq@hasabstractfalse
\long\def\banqAbstract#1#2{%
  \ifbundle
    \long\gdef\banq@abstracttext{#1}%
    \long\gdef\banq@keywordstext{#2}%
    \global\banq@hasabstracttrue
  \else
    \begin{abstract}#1\end{abstract}\par\medskip\noindent
    \textbf{Keywords:} #2\par
  \fi
}
\makeatother

% \banqPaper{<title>}{<subtitle>}
% Paper-level divider page printed at the top of each child's body when in
% bundle mode. Sits below the \part page in the visual hierarchy and provides
% a clear transition between consecutive papers within the same part (e.g.
% banq-lck -> banq-log inside Part~II). A no-op in standalone builds.
%
% Wrapped in \makeatletter so the macro body can reference \banq@RestoreLevels,
% \banq@DemoteLevels, and \@empty as single tokens rather than \banq + @… .
\makeatletter
\newcommand{\banqPaper}[2]{%
  \ifbundle
    % Force the cover page to land on an odd (recto) page; see \banqForceRecto.
    \banqForceRecto
    \onecolumn
    \thispagestyle{empty}
    % If a \banqPartWithEmoji call is pending, write its TOC + PDF-outline
    % entry now so the captured page number is the cover page (the visible
    % part boundary), not whatever blank pad the build happened to land on.
    \ifbanq@parttocpending
      \phantomsection
      \addcontentsline{toc}{part}{\thepart\hspace{1em}\banq@parttitle}%
      \global\banq@parttocpendingfalse
    \fi
    \vspace*{6em}
    \begin{center}
      {\Large\color{codegray} \textsc{XPower Banq: Part \Roman{part}\banqSubpartLetter}}\\[0.6em]
      \rule{0.4\textwidth}{0.4pt}\\[1.2em]
      {\fontsize{22}{26}\selectfont\bfseries #1}\\[0.8em]
      {\large\color{codegray} #2}
    \end{center}
    \vspace*{\fill}
    \ifbanq@hasemoji
      \begin{center}
        \twemoji[height=8cm]{\banq@partemoji}
      \end{center}
    \fi
    \vspace*{\fill}
    \clearpage
    % Verso page after the cover. If the child supplied an abstract via
    % \banqAbstract, render it here (still in onecolumn, unstyled body width
    % so the prose breathes). Otherwise leave the page blank — used for the
    % "References & Glossary" cover, which has no per-paper abstract since
    % the bundle's front-matter abstract already summarises the whole volume.
    \thispagestyle{empty}
    \ifbanq@hasabstract
      \vspace*{\fill}
      \begin{center}{\large\bfseries Abstract}\end{center}
      \medskip
      \begingroup
        \small
        \noindent\banq@abstracttext\par
        \medskip
        \noindent\textbf{Keywords:} \banq@keywordstext\par
      \endgroup
      \vspace*{\fill}
      \global\banq@hasabstractfalse
    \else
      \hbox{}%
    \fi
    \newpage
    \twocolumn
    % Register the divider in the TOC and PDF outline as a child of the part.
    % Restore section bookmark levels first so this divider always lands at
    % level 1 (regardless of any prior demotion left over from a sibling
    % sub-paper), then conditionally demote so subsequent sections nest under
    % this divider when we're inside a shared part.
    \banq@RestoreLevels
    \phantomsection
    \addcontentsline{toc}{section}{\protect\numberline{\Roman{part}\banqSubpartLetter.}#1}%
    \ifbanqSharedPart
      \banq@DemoteLevels
    \fi
  \fi
}
\makeatother


% ==============================================================================
% Title and Authors
% ==============================================================================
\title{\textbf{\Large XPower Banq: Mathematical Theory \& Proofs}\\[0.3em]
\normalsize Foundations, Formal Proofs, and Nash Equilibrium Analysis\\
for the XPower Banq Lending Protocol}

\author{
    \banqAuthor\thanks{\banqAuthorContact}
}

\date{}

\ifbundle\else
\renewcommand{\thepage}{MTP/\arabic{page}}
\fi

% ==============================================================================
% Document Begin
% ==============================================================================
\begin{document}

% Standalone-only front matter: title and arXiv tag (the abstract is emitted
% via \banqAbstract below, which renders it inline here in standalone mode
% and stashes it for the bundle's verso page after the cover).
\ifbundle\else
\maketitle
% ==============================================================================
% arxiv-tag.tex
% The arXiv-style vertical preprint tag printed on the first page of each
% standalone paper and on the bundle title page. The date is taken from
% \banqArxivDate (defined in banq-meta.tex) and may be overridden locally
% before % ==============================================================================
% arxiv-tag.tex
% The arXiv-style vertical preprint tag printed on the first page of each
% standalone paper and on the bundle title page. The date is taken from
% \banqArxivDate (defined in banq-meta.tex) and may be overridden locally
% before % ==============================================================================
% arxiv-tag.tex
% The arXiv-style vertical preprint tag printed on the first page of each
% standalone paper and on the bundle title page. The date is taken from
% \banqArxivDate (defined in banq-meta.tex) and may be overridden locally
% before \input{../shared/arxiv-tag} via \renewcommand.
% ==============================================================================

\begin{tikzpicture}[remember picture, overlay]
  \node[rotate=90, anchor=center, font=\Huge, color=codegray]
    at ([xshift=1cm]current page.west)
    {\banqArxivId~~\banqArxivDate};
\end{tikzpicture}
 via \renewcommand.
% ==============================================================================

\begin{tikzpicture}[remember picture, overlay]
  \node[rotate=90, anchor=center, font=\Huge, color=codegray]
    at ([xshift=1cm]current page.west)
    {\banqArxivId~~\banqArxivDate};
\end{tikzpicture}
 via \renewcommand.
% ==============================================================================

\begin{tikzpicture}[remember picture, overlay]
  \node[rotate=90, anchor=center, font=\Huge, color=codegray]
    at ([xshift=1cm]current page.west)
    {\banqArxivId~~\banqArxivDate};
\end{tikzpicture}

\fi

\banqAbstract{%
We collect the theoretical foundations of the XPower Banq lending protocol~\cite{wp}: continuous-compounding interest accrual, log-space TWAP construction, time-weighted parameter integration, token-bucket rate limiting, and formal proofs of cascade attenuation, Sybil rate-limiting, governance rate bounds, and MEV front-running resistance. We then analyse the Nash equilibrium of lock adoption under utilization-dependent rates and a secondary-market discount, characterising two self-reinforcing equilibria (low and high adoption) and proving that the protocol margin remains within sustainable bounds for any aggregate lock fraction $\bar{\rho} \in [0,1]$. Numerical evaluation of these results is the subject of the companion simulations paper~\cite{sim}.%
}{DeFi lending, formal proofs, Nash equilibrium, lock adoption, oracle theory, TWAP}

% ==============================================================================
% Body content
% ==============================================================================

\banqPaper{Mathematical Theory \& Proofs}{Foundations, Formal Proofs, and Nash Equilibrium Analysis for the XPower Banq Lending Protocol}

\section{Mathematical Foundations}
\label{apx:math}

\subsection{Interest Accrual}

Interest accrues via continuous compounding with ray precision (27 decimals).

\begin{definition}[Position Index]
The global index $I_t$ at time $t$ with annual rate $r$:
\begin{equation}
I_t = I_0 \cdot e^{r \cdot (t - t_0) / T_{\text{year}}}
\end{equation}
where $T_{\text{year}} = 365.25 \times 86400 = 31{,}557{,}600$ seconds.\footnote{\texttt{YEAR = CENTURY/100} with \texttt{CENTURY = 365\_25 days}, i.e.\ $36525 \cdot 86400 / 100 = 31{,}557{,}600$ s, to keep the constant integer-valued in Solidity.}
\end{definition}

\begin{theorem}[Balance Accrual]
A user's total balance $B_t$ given principal $P$ and checkpoint index $I_u$:
\begin{equation}
B_t = P \cdot \frac{I_t}{I_u}
\end{equation}
\end{theorem}

\begin{remark}[Piecewise integration in practice]
In the shipped contracts, $r$ is not a single constant but a function of utilization $r(t) = \hat r(\text{util}(t))$, and the index is integrated piecewise:
\begin{equation*}
\log I_t - \log I_0 = \sum_i \hat r(\text{util}_i) \cdot (t_{i+1} - t_i)/T_{\text{year}}
\end{equation*}
with re-indexings driven by every state-changing call. See \texttt{Position.\_indexOf} and \texttt{Position.\_reindex}; this matches the time-weighted framing of \S\ref{def:time-weighted-mean}.
\end{remark}

\subsection{Time-Weighted Integration}

\begin{definition}[Time-Weighted Mean]
\label{def:time-weighted-mean}
For parameter values $\theta_i$ at times $t_i$:
\begin{equation}
\bar{\theta}(t) = \frac{\sum_{i=0}^{n-1} \theta_i \cdot (t_{i+1} - t_i) + \theta_n \cdot (t - t_n)}{t - t_0}
\end{equation}
\end{definition}

\subsection{Oracle Price Aggregation}

\begin{definition}[Log-Space TWAP Quote]
\label{def:twap-quote}
Each refresh queries the AMM bidirectionally and computes:
\begin{align}
m_t &= \log_2(\bar{p}_t) && \text{(log mid price)} \\
s^{AB}_{t,log} &= \log_2(1 + s^{AB}_t) && \text{(log spread AB)} \\
s^{BA}_{t,log} &= \log_2(1 + s^{BA}_t) && \text{(log spread BA)} \\
r_t &= \tfrac{1}{2}\bigl(s^{AB}_{t,log} + s^{BA}_{t,log}\bigr) && \text{(log geomean spread)}
\end{align}
The EMA smooths in log space with decay $\alpha$, blending the previous mean with the \emph{previously stored} sample (the contract maintains a one-sample buffer: \texttt{TWAPLib.update} reads \texttt{last\_mid = midOf(last)} \emph{before} overwriting \texttt{twap\_.last} with the new $m_t$, so the new sample enters the mean only at the next refresh):
\begin{align}
\overline{m}_t &= \alpha \cdot \overline{m}_{t-1} + (1{-}\alpha) \cdot m_{t-1} \\
\overline{r}_t &= \alpha \cdot \overline{r}_{t-1} + (1{-}\alpha) \cdot r_{t-1}
\end{align}
where $m_{t-1}$ (resp.\ $r_{t-1}$) denotes the most recent sample observed strictly \emph{before} the current refresh.
\end{definition}

\begin{definition}[Spread Scaling]
\label{def:spread-scaling}
Let $n = \texttt{amount}/\texttt{unit\_source}$ be the position expressed in \emph{whole units} of the source asset (so a 1.5 ETH order has $n = 1.5$, not $1.5 \cdot 10^{18}$). With base spread $s = 2^{\overline{r}} - 1$:
\begin{align}
x &= n \cdot s, \quad \mu = \log_2(2x + 2), \quad s' = s \cdot \mu
\end{align}
Final quotes: $\text{bid}(n) = \text{value} - s' \cdot \text{value}/2$, $\text{ask}(n) = \text{value} + s' \cdot \text{value}/2$. Implemented in \texttt{Oracle.\_hlfSpread} / \texttt{Oracle.\_center}.
\end{definition}

\subsection{Rate Limiting}

\begin{definition}[Token Bucket]
A rate limiter with capacity $C_{\max}$, regeneration rate $1$ unit per wall-clock second (a hard-coded design choice, not a governable parameter), and per-call cost $c$ supplied by the caller (the \texttt{floor\_cost} argument of each modifier invocation):
\begin{equation}
C'(t) = \min\left(C_{\max}, C(t_{\text{last}}) + (t - t_{\text{last}})\right) - c
\end{equation}
Operation allowed iff $C'(t) \geq 0$. Implementation: \texttt{RateLimited.\_ratelimitedCheck} regenerates first, reverts if $C(t_{\text{last}}) + (t - t_{\text{last}}) < c$, then commits the post-decrement state---algebraically equivalent to the post-cost residual $C'(t) \geq 0$.
\end{definition}

\section{Formal Proofs}
\label{apx:proofs}

\subsection{Cascade Attenuation Proof}

\begin{proof}[Proof of the Cascade Attenuation Theorem (\S4.1 of~\cite{wp})]
Let liquidation of position $P$ with lock fraction $\phi$ transfer supply tokens $S$ to liquidator $L$.

\textbf{Step 1:} By \texttt{Position.burn}, \texttt{redeem} reverts when the post-burn balance would fall below the user's aggregate lock total: the post-condition asserts \texttt{totalOf(user) $\geq$ \_lock.totalOf(user)} after burning the requested amount, raising \texttt{Locked(user, total)} otherwise. There is no per-token ``locked vs.\ unlocked'' tag; the lock is a reservation against the aggregate balance. Hence the locked share $\phi \cdot S$ of any position cannot leave the protocol; only the unlocked surplus $(1{-}\phi) \cdot S$ is available for redemption and resale.

\textbf{Step 2:} Liquidator $L$ receives position tokens, not underlying assets. The underlying assets remain in the Vault.

\textbf{Step 3:} Only unlocked fraction $(1{-}\phi)$ can generate sell pressure. Price impact: $\Delta p = f((1{-}\phi) \cdot V_{\text{sold}})$.

\textbf{Step 4:} Lock property is preserved through transfer: by \texttt{Position.\_lockPush} $\to$ \texttt{LockLib.push}, each ring-buffer slot's value is scaled by $\texttt{mulDiv}(\texttt{src\_value}, \texttt{amount}, \texttt{balance})$ and added to the target. This is invoked unconditionally on every \texttt{transferFrom}, including the liquidation path. Hence the per-user lock fraction $\phi_{\text{user}} = \texttt{depthOf}(\text{user})/\texttt{balance}(\text{user})$ is preserved across any transfer of size $\leq$ \texttt{balance}. Therefore the cascade amplification factor is bounded by $(1{-}\phi)$.
\end{proof}

\subsection{Sybil Rate-Limiting}

\begin{definition}[Sybil Resistance Scope]
\label{def:sybil-scope}
The cap system bounds accumulation \emph{rate} via holder-count scaling. It does \textbf{not} guarantee that more accounts hurts an attacker's long-term equilibrium share---simulation shows $<$2\% penalty over 60 weeks.
\end{definition}

\begin{theorem}[Sybil Attack Cost]
\label{thm:sybil-cost}
Assume each of the $k$ Sybil accounts holds a positive balance of at least one whole position token (\texttt{balance $\geq$ unit\_token = $10^{\text{decimals}}$}), so that it is counted as a \emph{large holder} by \texttt{Position.\_largeHolders} (the running counter incremented in \texttt{\_update} only when an account crosses \texttt{balance $\geq$ unit}). Assume further that the per-account cap considered is for a \emph{subsequent} deposit on top of an existing positive balance, so that the $\beta(\lambda) = 12\lambda(1{-}\lambda)^2$ weighting of \texttt{\_capOf} (gated on \texttt{balance > 0 \&\& total > balance}) applies. Then for the attacker to increase total cap gain by factor $F$, they must create $k > F^2 \cdot (n{+}2) - (n{+}2)$ such large-holder Sybils. This quadratic scaling makes large gains expensive.
\end{theorem}

\begin{proof}
Under the large-holder hypothesis the denominator $\sqrt{\texttt{largeHolders()}+2}$ scales from $\sqrt{n+2}$ to $\sqrt{n+k+2}$ when $k$ funded Sybils are added; a Sybil account with \texttt{balance < unit} would not advance the counter and is excluded by hypothesis. Under the existing-holder hypothesis, the beta block of \texttt{\_capOf} is active. With 1 account: $C_1 = C_{\max} \cdot \beta(\lambda) / \sqrt{n{+}2}$. With $k$ equal accounts and small $\lambda/k$: $C_k \approx 12\lambda C_{\max}/\sqrt{n{+}k{+}2}$. For $C_k > F \cdot C_1$: $\sqrt{n{+}2}/\sqrt{n{+}k{+}2} > F/k$, yielding the stated bound. (A floor \texttt{MIN\_HOLDERS\_ID} on \texttt{largeHolders()} prevents the denominator from collapsing when $n$ is small.)
\end{proof}

\subsection{Governance Rate Bound Proof}

\begin{proof}[Proof of the Governance Rate Bound (\S6 of~\cite{wp})]
\textbf{Step 1 (Single bound):} Each transition satisfies $\theta_0/\mu \leq \theta_1 \leq \mu \cdot \theta_0$, enforced inline by \texttt{Parameterized.\_setTargetIf}, which compares the proposed value against \texttt{old\_value <{}< 1} (upper bound, $\mu = 2$) and \texttt{old\_value >{}> 1} (lower bound, $1/\mu$) and reverts on violation.

\textbf{Step 2 (Rate limiting):} \texttt{setTarget} is guarded by the \texttt{limited(Constant.MONTH, \dots)} modifier, which restricts updates to once per $\Delta t_{\min} = $ \texttt{Constant.MONTH} $= $ \texttt{Constant.YEAR}$/12 \approx 30.4375$ days (a compile-time constant, not a governable parameter). In interval $(t{-}t_0)$, at most $n = \lfloor(t{-}t_0)/\Delta t_{\min}\rfloor$ transitions occur.

\textbf{Step 3 (Composition):} After $n$ transitions: $\theta_n \leq \mu^n \theta_0$ and $\theta_n \geq \mu^{-n} \theta_0$.

\textbf{Step 4 (Transitions):} During transitions, the time-weighted mean stays between $\theta_{\text{old}}$ and $\theta_{\text{new}}$, preserving the bound.
\end{proof}

\begin{corollary}[Attack Detection Window]
To change a parameter by factor $F > \mu$: $T_{\text{attack}} \geq \lceil\log_\mu F\rceil \cdot \Delta t_{\min}$. With $\mu = 2$ and the compile-time constant $\Delta t_{\min} = $ \texttt{Constant.MONTH} $\approx 30.4375$ days (not itself governable), a $10\times$ change requires 4~months.
\end{corollary}

\subsection{MEV Front-Running Resistance}

\begin{theorem}[MEV Front-Running Cost]
\label{thm:mev}
A liquidation with profit $\pi$ cannot be profitably front-run if:
$H_{\text{attacker}}/H_{\text{victim}} < t_{\text{cache}}/T_{\text{PoW}}(\pi)$
where $t_{\text{cache}}$ is the \texttt{PowLimited} cache window over which the PoW key is stable.
\end{theorem}

\begin{proof}
The PoW hash includes \texttt{tx.origin}, preventing solution reuse across attacker EOAs. The \texttt{PowLimited} modifier caches the block hash anchoring the PoW key for \texttt{cacheTime} seconds, refreshing only once \texttt{cacheTime + \_blockTime < block.timestamp}. Both \texttt{Pool} and \texttt{Oracle} instantiate \texttt{PowLimited(1 hours)}, so $t_{\text{cache}} = 1$~hour~$= 3600$~s. With expected solve time $T_{\text{PoW}}(\pi) = 16^d / H$, the attacker must mine independently within $t_{\text{cache}}$ against its own \texttt{tx.origin}; front-running fails when $T_{\text{PoW}} > t_{\text{cache}}$, i.e.\ when difficulty $d$ is tuned so that $16^d / H_{\text{attacker}} > t_{\text{cache}} = 3600$~s.
\end{proof}

\section{Nash Equilibrium Analysis for Lock Adoption}
\label{apx:nash}

This section analyzes the game-theoretic equilibrium of lock adoption. We model rational actors choosing between locked and unlocked positions, accounting for APY differentials, secondary market dynamics, and liquidation seniority.

\subsection{Model Setup}

\subsubsection{Position Types and Rates}

With spread $s = 10\%$ and $r_{\text{bonus}} = r_{\text{malus}} = s$:

\begin{center}
\small
\begin{tabularx}{0.95\columnwidth}{@{}l r >{\raggedleft\arraybackslash}X@{}}
\toprule
Position Type & Rate Multiplier & APY vs.\ Base \\
\midrule
Unlocked Supplier & $(1-s) = 0.90$ & $-10\%$ \\
Locked Supplier & $(1+r_{\text{bonus}})(1-s) = 0.99$ & $-1\%$ \\
Unlocked Borrower & $(1+s) = 1.10$ & $+10\%$ \\
Locked Borrower & $(1-r_{\text{malus}})(1+s) = 0.99$ & $-1\%$ \\
\bottomrule
\end{tabularx}
\end{center}

The \textbf{APY differential} between locked and unlocked positions:
\begin{equation}
\Delta r = r_{\text{base}} \times (1-s) \times r_{\text{bonus}} = r_{\text{base}} \times 0.9 \times 0.1 = 0.09 \times r_{\text{base}}
\end{equation}

\subsubsection{Secondary Market Dynamics}

Position tokens are ERC20-transferable, enabling secondary market trading. However, locked positions trade at a \textbf{discount} $D$ relative to NAV because:
\begin{enumerate}
    \item Locked positions cannot be redeemed for underlying assets
    \item Exit is only possible via secondary market sale
    \item Lock status transfers proportionally with tokens (buyers receive locked tokens)
\end{enumerate}

The discount $D$ is endogenous to market conditions, typically ranging from 3--10\%\footnote{Empirical range observed in secondary markets for illiquid DeFi positions: Lido stETH traded at 6--7\% discount during the June 2022 liquidity crisis (Nansen, ``On-Chain Forensics: Demystifying stETH's De-peg,'' 2022); Convex cvxCRV has traded at 5--12\% discount to CRV due to permanent lock-up (geeogi, ``cvxCRV peg,'' 2022).} based on:
\begin{itemize}
    \item Secondary market liquidity depth
    \item Expected holding period of market participants
    \item Overall lock adoption rate (affecting liquidity)
\end{itemize}

\subsection{Player Utility Functions}

\subsubsection{Supplier Utility}

For a supplier with lock ratio $\rho \in [0,1]$ and holding period $T$:

\textbf{Unlocked position} ($\rho = 0$):
\begin{equation}
U_{\text{unlocked}} = P \times (1 + r_{\text{unlocked}})^T
\end{equation}

\textbf{Locked position} ($\rho = 1$):
\begin{equation}
U_{\text{locked}} = P \times (1 + r_{\text{locked}})^T \times (1 - D)
\end{equation}

where the $(1-D)$ factor reflects the secondary market discount on exit.

\subsubsection{Breakeven Condition}

Lock adoption is rational when $U_{\text{locked}} > U_{\text{unlocked}}$:
\begin{equation}
(1 + r_{\text{locked}})^T \times (1 - D) > (1 + r_{\text{unlocked}})^T
\end{equation}

Solving for the \textbf{breakeven holding period} $T^*$:
\begin{equation}
\label{eq:breakeven}
T^* = \frac{\ln(1-D)}{\ln(1 + r_{\text{unlocked}}) - \ln(1 + r_{\text{locked}})} \approx \frac{D}{\Delta r}
\end{equation}

\subsection{Breakeven Analysis by Utilization}

The base interest rate $r_{\text{base}}$ varies with utilization $U$ according to the piecewise-linear model (\S4.7 of~\cite{wp}). Table~\ref{tab:nash-breakeven} presents breakeven periods for various utilization levels.

\begin{table}[htbp]
\centering
\caption{Breakeven holding periods by utilization}
\label{tab:nash-breakeven}
\small
\begin{tabularx}{0.95\columnwidth}{@{}>{\raggedleft\arraybackslash}X >{\raggedleft\arraybackslash}X >{\raggedleft\arraybackslash}X >{\raggedleft\arraybackslash}X >{\raggedleft\arraybackslash}X@{}}
\toprule
Utilization $U$ & Base Rate $r$ & Differential $\Delta r$ & $T^*$ ($D=5\%$) & $T^*$ ($D=3\%$) \\
\midrule
80\% & 8.9\% & 0.80\% & 6.3y & 3.8y \\
90\% & 10\% & 0.90\% & 5.6y & 3.3y \\
95\% & 55\% & 4.95\% & 12m & 7m \\
98\% & 82\% & 7.38\% & 8m & 5m \\
100\% & 100\% & 9.0\% & 6.7m & 4m \\
\bottomrule
\end{tabularx}
\end{table}

Lock adoption thus becomes economically attractive primarily during high-utilization periods ($U > 95\%$), precisely when cascade prevention is most valuable.

\subsection{Liquidation Seniority Value}

Locked positions gain \emph{de facto} seniority in liquidations because liquidators prefer unlocked positions (immediate liquidity). The seniority value $S(\bar{\rho})$ depends on aggregate lock adoption $\bar{\rho}$ (denoted $\phi$ in the Cascade Attenuation Theorem (\S4.1 of~\cite{wp})):

\begin{equation}
S(\bar{\rho}) \approx P(\text{liquidation}) \times (1 - \bar{\rho}) \times L_{\text{loss}}
\end{equation}

where $P(\text{liquidation})$ is the probability of a liquidation event and $L_{\text{loss}}$ is the expected loss in liquidation.

This creates a \textbf{coordination game}: seniority value decreases as lock adoption increases. At $\bar{\rho} = 0$, locked positions have maximum seniority; at $\bar{\rho} = 1$, no differentiation exists.

\subsection{Nash Equilibrium Characterization}

\subsubsection{Equilibrium Condition}

At equilibrium lock adoption $\bar{\rho}^*$, the marginal user is indifferent between locking and not locking:
\begin{equation}
\underbrace{\Delta r \times T}_{\text{APY benefit}} + \underbrace{S(\bar{\rho}^*)}_{\text{Seniority}} = \underbrace{D}_{\text{Discount cost}}
\end{equation}

\subsubsection{Utilization-Dependent Equilibria}

\begin{table}[htbp]
\centering
\caption{Equilibrium lock adoption by utilization regime}
\label{tab:nash-equilibria}
\small
\begin{tabularx}{0.95\columnwidth}{@{}l >{\raggedleft\arraybackslash}X >{\raggedleft\arraybackslash}X >{\raggedleft\arraybackslash}X@{}}
\toprule
Utilization Regime & Base Rate & Expected $\bar{\rho}^*$ & Protocol Margin \\
\midrule
Normal ($U \in [0,90\%)$) & $<10\%$ & 10--20\% & 16--18\% of $r$ \\
Elevated ($U \in [90\%,95\%)$) & 10--55\% & 20--40\% & 12--16\% of $r$ \\
Stressed ($U \in [95\%,100\%]$) & $>55\%$ & 40--70\% & 6--12\% of $r$ \\
\bottomrule
\end{tabularx}
\end{table}

The mechanism thus aligns incentives with protocol needs: under normal conditions, low lock adoption preserves revenue, while under stressed conditions, high adoption provides cascade protection when it is most needed.

\subsection{Protocol Margin Analysis}

\subsubsection{Margin as Function of Lock Adoption}

Protocol margin per unit interest flow:
\begin{equation}
M(\bar{\rho}) = 2s - r_{\text{bonus}} \times \bar{\rho}^S - r_{\text{malus}} \times \bar{\rho}^B
\end{equation}

With $r_{\text{bonus}} = r_{\text{malus}} = s$ and assuming $\bar{\rho}^S = \bar{\rho}^B = \bar{\rho}$:
\begin{equation}
M(\bar{\rho}) = 2s(1 - \bar{\rho})
\end{equation}

\begin{table}[htbp]
\centering
\caption{Protocol margin by lock adoption}
\label{tab:nash-margin}
\small
\begin{tabularx}{0.95\columnwidth}{@{}>{\raggedleft\arraybackslash}X >{\raggedleft\arraybackslash}X >{\raggedleft\arraybackslash}X@{}}
\toprule
Lock Adoption $\bar{\rho}$ & Margin & Retention \\
\midrule
0\% & 20\% of $r$ & 100\% \\
25\% & 15\% of $r$ & 75\% \\
50\% & 10\% of $r$ & 50\% \\
75\% & 5\% of $r$ & 25\% \\
100\% & 0\% of $r$ & 0\% \\
\bottomrule
\end{tabularx}
\end{table}

\subsubsection{Solvency at Full Adoption}

The default configuration operates at the solvency boundary:
\begin{equation}
r_{\text{bonus}} + r_{\text{malus}} = s + s = 2s \quad \checkmark
\end{equation}

Protocol solvency is maintained for any $\bar{\rho} \in [0, 1]$, though margin approaches zero as $\bar{\rho} \to 1$.

\subsection{Stability Analysis}

\subsubsection{Multiple Equilibria}

The lock adoption game exhibits \textbf{multiple equilibria} due to the endogenous discount:

\begin{enumerate}
    \item \textbf{Low-adoption equilibrium} ($\bar{\rho}^* \approx 10\%$): Thin secondary market $\Rightarrow$ high discount $D$ $\Rightarrow$ locking unattractive $\Rightarrow$ low adoption (self-reinforcing)

    \item \textbf{High-adoption equilibrium} ($\bar{\rho}^* \approx 60\%$): Deep secondary market $\Rightarrow$ low discount $D$ $\Rightarrow$ locking attractive $\Rightarrow$ high adoption (self-reinforcing)
\end{enumerate}

\subsubsection{Equilibrium Selection}

The realized equilibrium depends on:
\begin{itemize}
    \item \textbf{Initial conditions}: Early lock adoption seeds secondary market liquidity
    \item \textbf{Utilization dynamics}: High-utilization periods push toward high-adoption equilibrium
    \item \textbf{External liquidity}: Protocol-seeded liquidity can bootstrap the high-adoption equilibrium
\end{itemize}

\subsection{Summary}

The lock mechanism creates a self-regulating system in which rational actors provide cascade protection precisely when the protocol most needs it.

\begin{theorem}[Lock Adoption Equilibrium]
\label{thm:nash}
With $r_{\text{bonus}} = r_{\text{malus}} = s = 10\%$:
\begin{enumerate}
    \item Lock adoption is \textbf{utilization-dependent}: attractive when $U > 95\%$, marginal otherwise
    \item The APY differential $\Delta r = 9\% \times r_{\text{base}}$ provides maximum incentive within solvency constraints
    \item Lock adoption \textbf{self-regulates}: increases during stress (cascade protection) and decreases during calm (revenue preservation)
    \item Protocol margin ranges from 20\% of $r$ (no locks) to 0\% of $r$ (full locks), with expected operating range 8--16\% of $r$
\end{enumerate}
\end{theorem}

The secondary market discount acts as a natural governor---preventing full adoption while ensuring sufficient incentive for meaningful cascade protection---so the mechanism achieves incentive alignment without external intervention.

The welfare balance is favorable. Cascade attenuation reduces depth by factor~$(1{-}\phi)$, and spread compression rewards committed participants---at full lock with $s = r_{\text{bonus}} = r_{\text{malus}} = 10\%$, both locked suppliers and borrowers face effective rate multiplier~$0.99$, approaching a zero-spread market. Against these benefits stands liquidity fragmentation: locked supply raises effective utilization to $U_{\text{eff}} = V_b / (V_s(1{-}\phi))$, potentially triggering kink-rate dynamics at lower borrow volumes. The interest rate curve absorbs this naturally---higher effective utilization raises rates, incentivizing settlement---so the effect narrows the sub-kink operating range without creating instability.

Together, these properties confirm Theorem~\ref{thm:nash}: the mechanism is self-regulating, with per-parameter solvency bounds that require no on-chain adoption tracking, cascade dynamics that naturally throttle liquidation velocity, and welfare effects that the interest rate curve absorbs without instability.

\subsubsection{Equilibrium Structure}

The game exhibits two self-reinforcing equilibria: high adoption ($\bar{\rho} > 40\%$, margin 8--12\% of~$r$, strong cascade protection) and low adoption ($\bar{\rho} < 15\%$, margin 16--20\% of~$r$, high cascade exposure). Selection is path-dependent---seeding secondary market liquidity bootstraps the high-adoption regime, while neglect traps the protocol in low adoption. Governance can influence this by initially setting $r_{\text{bonus}}$ above the long-run target, then reducing it via lethargic transition once adoption stabilizes.

Within either regime, the mechanism is countercyclical: lock adoption rises above $U = 95\%$ when cascade prevention is most valuable, then recedes as utilization normalizes, restoring margin. Lethargic governance reinforces this---the 0.5x--2x per-cycle parameter bound ensures that equilibrium shifts occur gradually enough for the secondary market to adjust. During stress, protocol margin decreases as the intended trade-off for cascade protection, and recovers during the subsequent normalization.

The equilibrium is asymmetric across position types. Locked suppliers earn at effective rate $r_{\text{base}}(1{-}s)(1{+}r_{\text{bonus}} \cdot \rho_u)$ but must compensate the secondary market discount~$D$, requiring $U > 95\%$ for rational adoption. Locked borrowers pay at effective rate $r_{\text{base}}(1{+}s)(1{-}r_{\text{malus}} \cdot \rho_u)$---reducing an \emph{obligation} with no redemption friction---making $\bar{\rho}^B > \bar{\rho}^S$ in the normal regime, converging during stress as the supply bonus dominates. A plausible operating point ($\bar{\rho}^S \approx 15\%$, $\bar{\rho}^B \approx 35\%$) yields:
\begin{equation}
M = 2s - s \cdot \bar{\rho}^S - s \cdot \bar{\rho}^B = 20\% - 1.5\% - 3.5\% = 15\% \text{ of } r
\end{equation}
within the sustainable range of Table~\ref{tab:nash-margin}.

\subsubsection{Cascade Dynamics and Keeper Sizing}

Debt assumption transfers locks proportionally, preserving aggregate~$\bar{\rho}$ through cascades---the attenuation bound $(1{-}\phi)$ from Theorem~4.1 of~\cite{wp} holds dynamically, not merely at a pre-crash snapshot. At $\bar{\rho} \approx 30\%$, this reduces 25\%-shock liquidations from 85.7\% to~$\approx$55\% of positions (Table~5 of~\cite{wp}). However, locked supply received by a liquidator cannot be redeemed: absorbing a fully-locked victim at exponent~$e$ increases assumed debt by $d = V_b \cdot 2^{-e}$ while adding zero redeemable supply, eroding headroom for subsequent liquidations.

The exponent choice is itself strategic. Choosing $e = 0$ captures the full position---maximizing per-event profit but inheriting the full lock fraction. Choosing $e \geq 1$ halves the position taken, reducing lock inheritance at the cost of leaving the remainder available to competitors. In equilibrium, the optimal~$e$ reflects each liquidator's current health factor and the remaining cascade depth: unconstrained liquidators with ample headroom prefer small~$e$, while capacity-constrained ones ration across targets. Both effects---lock removal from the liquidation pool and liquidator self-throttling---reinforce cascade attenuation.

Keeper sizing must account for lock propagation. A keeper processing $k$ sequential liquidations needs headroom for both aggregate debt absorbed and the locked fraction that cannot be redeemed. If the expected victim lock fraction is~$\bar{\rho}$, the effective headroom requirement per liquidation increases by factor $1/(1{-}\bar{\rho})$ relative to the unlocked case, since only $(1{-}\bar{\rho})$ of received supply contributes redeemable collateral. The cascade simulation in the companion simulations paper~\cite{sim} quantifies this across varying lock fractions and crash severities.

\subsubsection{Solvency Guarantees}

Protocol solvency is enforced algebraically. The supervised contract enforces $r_{\text{bonus}} \leq s$ and $r_{\text{malus}} \leq s$ in \emph{both} directions via \texttt{\_setTarget} overrides: when setting \texttt{LOCK\_BONUS\_ID}/\texttt{LOCK\_MALUS\_ID}, the new value must be $\leq s$; when setting \texttt{SPREAD\_ID}, the new spread must remain $\geq \max(r_{\text{bonus}}, r_{\text{malus}})$, so a governance-driven \emph{decrease} of $s$ cannot violate the invariant either. Additionally $s \leq \texttt{Constant.HLF} = 0.5$ is enforced. Since $\bar{\rho}^S, \bar{\rho}^B \in [0,1]$, the aggregate condition $r_{\text{bonus}} \bar{\rho}^S + r_{\text{malus}} \bar{\rho}^B \leq 2s$ holds without on-chain lock-adoption tracking. The worst case ($\bar{\rho}^S = \bar{\rho}^B = 1$) yields zero margin but never insolvency---avoiding the gas overhead of global lock-adoption statistics by relying on per-parameter bounds enforced at governance time. Lock adoption also tightens beta-distributed caps indirectly: locked suppliers are persistent holders, inflating the effective~$n$ in the $\sqrt{n{+}2}$ divisor, while attackers splitting across $k$ accounts face the $O(\sqrt{k})$ cap-scaling barrier \emph{and} irrevocability---acquiring locked positions via secondary market incurs discount~$D$, creating a multiplicative barrier to rapid concentration.

\subsubsection{Welfare}

The welfare balance is favorable. Cascade attenuation reduces depth by factor~$(1{-}\phi)$, and spread compression rewards committed participants---at full lock with $s = r_{\text{bonus}} = r_{\text{malus}} = 10\%$, both locked suppliers and borrowers face effective rate multiplier~$0.99$, approaching a zero-spread market. Against these benefits stands liquidity fragmentation: locked supply raises effective utilization to $U_{\text{eff}} = V_b / (V_s(1{-}\phi))$, potentially triggering kink-rate dynamics at lower borrow volumes. The interest rate curve absorbs this naturally---higher effective utilization raises rates, incentivizing settlement---so the effect narrows the sub-kink operating range without creating instability.

Together, these properties confirm Theorem~\ref{thm:nash}: the mechanism is self-regulating, with per-parameter solvency bounds that require no on-chain adoption tracking, cascade dynamics that naturally throttle liquidation velocity, and welfare effects that the interest rate curve absorbs without instability.

% ==============================================================================
% Standalone-only References
% ==============================================================================
\ifbundle\else
\begin{thebibliography}{99}

\bibitem{wp}
\textsc{Karun The Ritch}.
\newblock XPower Banq: A Permissionless DeFi Lending Protocol with Lethargic Governance and Beta-Distributed Position Caps.
\newblock \emph{Preprint}, 2026.
\newblock \url{https://www.xpowerbanq.com}

\bibitem{sim}
\textsc{Karun The Ritch}.
\newblock XPower Banq: Simulations \& Risk Analysis.
\newblock \emph{Preprint}, 2026.
\newblock \url{https://www.xpowerbanq.com}

\bibitem{ref}
\textsc{Karun The Ritch}.
\newblock XPower Banq: References \& Glossary.
\newblock \emph{Preprint}, 2026.
\newblock \url{https://www.xpowerbanq.com}

\end{thebibliography}
\fi

% ==============================================================================
% Glossary
% ==============================================================================
\ifbundle\else
\section*{Glossary}
\label{mtp:glossary}
\addcontentsline{toc}{section}{Glossary}

\noindent A comprehensive glossary with over 120 entries appears in \cite{ref}. The following terms are central to this paper.

\begin{description}[style=nextline, font=\normalfont\bfseries, leftmargin=1.5em]

\item[$\alpha$\,$\dagger$] EMA decay parameter of the log-space TWAP oracle. Governs tracking speed: $\alpha = 0.5^{1/\text{HL}}$ where HL is the half-life in refreshes. Default $\alpha \approx 0.944$ (12 refreshes; with the default 1-hour refresh interval enforced by \texttt{PowLimited(1 hours)} in \texttt{Oracle.sol}, this is a 12-hour half-life).

\item[Asymptotic Transition] Parameter change mechanism where values approach targets gradually via time-weighted mean rather than discrete jumps. See Lethargic Governance.

\item[Beta Cap] Position limit following \mbox{$12\lambda(1-\lambda)^2$} distribution where $\lambda$ is the user's fraction of total supply; peaks at $\lambda = 1/3$ and vanishes at boundaries.

\item[Cascade Attenuation] Reduction of liquidation-cascade depth by factor~$(1{-}\phi)$, where $\phi$ is the locked fraction of the position; the central result of \S4.1 of~\cite{wp}.

\item[Coordination Game] Strategic interaction where player payoffs depend on aggregate choices; lock adoption exhibits coordination dynamics where seniority value decreases as adoption increases.

\item[Geomean Spread] Bidirectional geometric mean of relative spreads from forward and reverse AMM queries; stored in log-space as $\log_2(1 + s_{\text{geo}})$. Provides symmetric, manipulation-resistant spread estimation.

\item[Half-Life] Time for EMA weight to decay to 50\%; determines TWAP responsiveness to new price observations. \cite[\S C]{wp} tabulates decay factors.

\item[Health Factor] Ratio of weighted supply value to weighted borrow value: $H = \sum w_s V_s / \sum w_b V_b$. Liquidation occurs when $H < 1$.

\item[Lethargic Governance] Governance model with time-delayed, bounded parameter transitions. Values approach targets asymptotically, bounded to $0.5\times$--$2\times$ per governance cycle. See \cite[\S A]{wp}.

\item[Liquidation Cascade] Destructive feedback loop where forced sales depress prices, triggering further liquidations. Modeled in \cite{sim}.

\item[Lock Adoption] Aggregate fraction $\bar{\rho}$ of positions that are locked across the protocol; equilibrium adoption varies with utilization regime. Analyzed in this paper.

\item[Log-Space Oracle] Oracle architecture storing prices as $\log_2(\text{price})$ and spreads as $\log_2(1 + s)$; enables EMA smoothing as geometric-mean temporal averaging.

\item[Multiplicative Bounds] Constraint limiting parameter changes to $0.5\times$--$2\times$ per governance cycle; prevents rapid manipulation.

\item[Nash Equilibrium] Stable strategic state where no player can improve their payoff by unilaterally changing strategy; lock adoption exhibits utilization-dependent equilibria. Analyzed in this paper.

\item[Protocol Margin] Revenue retained by protocol from interest spread; $M(\bar{\rho}) = 2s(1-\bar{\rho})$ where $s$ is spread and $\bar{\rho}$ is lock adoption rate.

\item[Solvency Boundary] Parameter constraint ensuring protocol can meet all obligations; default configuration satisfies $r_{\text{bonus}} + r_{\text{malus}} = 2s$, maintaining solvency for any lock adoption rate.

\item[Spread] (1)~Oracle: bidirectional geomean of rel.\ spreads from both AMM query directions, stored as $\log_2(1 + s_{\text{geo}})$; wider spreads indicate lower liquidity or higher manipulation resistance. (2)~IRM: symmetric half-spread parameter (e.g., 10\%) applied to base rate; borrow rate = base~$\times$~$(1{+}s)$, supply rate = base~$\times$~$(1{-}s)$.

\item[Time-Weighted Mean] Integration technique averaging parameter values over time; enables smooth transitions without discrete jumps.

\item[Token Bucket] Rate limiting mechanism with capacity $C$, regeneration rate, and per-operation cost; operation allowed iff $C \geq 0$.

\item[Utilization] Ratio of borrowed assets to supplied assets in a vault; drives interest rates via a piecewise-linear model with a kink at optimal utilization.

\end{description}

\clearpage
\ifodd\value{page}\else\thispagestyle{empty}\null\clearpage\fi
\fi

\end{document}
